% ============================================================================
% LANGENTWURF: Pretérito Indefinido (Einführung -ar Verben)
% ============================================================================
% Sequenz: 08 - Las vacaciones
% Stunde: 08c - Pretérito Indefinido Einführung
% Dauer: 90 Min. (Doppelstunde)
% ============================================================================

\documentclass[11pt,a4paper]{article}

% ============================================================================
% PACKAGES
% ============================================================================
\usepackage[utf8]{inputenc}
\usepackage[T1]{fontenc}
\usepackage[ngerman]{babel}
\usepackage{geometry}
\usepackage{fancyhdr}
\usepackage{lastpage}
\usepackage{xcolor}
\usepackage{tcolorbox}
\usepackage{enumitem}
\usepackage{tabularx}
\usepackage{longtable}
\usepackage{array}
\usepackage{booktabs}
\usepackage{hyperref}
\usepackage{ifthen}
\usepackage{amssymb}

% ============================================================================
% KONFIGURATION
% ============================================================================

\newcommand{\plantier}{1}
\newcommand{\fachid}{spanisch}
\newcommand{\fach}{Spanisch}
\newcommand{\fachfarbe}{c41e3a}

% ============================================================================
% VARIABLEN
% ============================================================================

\newcommand{\jahrgangsstufe}{7}
\newcommand{\schulart}{Gesamtschule}
\newcommand{\stundenthema}{Pretérito Indefinido (Einführung -ar)}
\newcommand{\reihenthema}{Las vacaciones}
\newcommand{\stundennummer}{8c}
\newcommand{\reihenstunden}{4}
\newcommand{\unterrichtsdatum}{\underline{\hspace{3cm}}}
\newcommand{\unterrichtszeit}{\underline{\hspace{3cm}}}
\newcommand{\raum}{\underline{\hspace{2cm}}}

\newcommand{\lehrkraft}{N.N.}
\newcommand{\ausbildungsstand}{Lehrkraft}

\newcommand{\lerngruppengroesse}{24}
\newcommand{\maennlich}{12}
\newcommand{\weiblich}{12}
\newcommand{\divers}{0}

\newcommand{\gerniveau}{A1}
\newcommand{\lehrwerk}{Qué pasa 1}
\newcommand{\lehrwerkseite}{S. 94--97}

% ============================================================================
% LAYOUT
% ============================================================================
\geometry{left=2.5cm, right=2.5cm, top=2.5cm, bottom=2.5cm}
\definecolor{fach}{HTML}{\fachfarbe}
\definecolor{gray}{HTML}{666666}
\definecolor{lightgray}{HTML}{f5f5f5}
\definecolor{successgreen}{HTML}{2a7a4b}
\definecolor{warningorange}{HTML}{b35c00}

\tcbuselibrary{skins,breakable}

\newtcolorbox{infobox}[1][]{
    colback=lightgray,
    colframe=fach,
    fonttitle=\bfseries,
    title=#1,
    breakable
}

\newtcolorbox{scorebox}{
    colback=white,
    colframe=fach,
    fonttitle=\bfseries,
    title=Didaktik-Score,
    breakable
}

\pagestyle{fancy}
\fancyhf{}
\fancyhead[L]{\small\textcolor{gray}{\fach{} -- Jg. \jahrgangsstufe{} -- \stundenthema}}
\fancyhead[R]{\small\textcolor{gray}{Langentwurf Tier 1}}
\fancyfoot[C]{\small\thepage{} / \pageref{LastPage}}
\renewcommand{\headrulewidth}{0.4pt}
\renewcommand{\footrulewidth}{0pt}

% ============================================================================
% DOKUMENT
% ============================================================================
\begin{document}

% ============================================================================
% DECKBLATT
% ============================================================================
\begin{titlepage}
    \centering
    \vspace*{2cm}

    {\large\textcolor{gray}{\schulart}}\\[2cm]

    {\Huge\textcolor{fach}{\textbf{Langentwurf}}}\\[0.5cm]
    {\large Tier 1 -- Vollständig}\\[2cm]

    {\LARGE\textbf{\stundenthema}}\\[0.5cm]
    {\large im Rahmen der Unterrichtsreihe}\\[0.3cm]
    {\Large\textit{\reihenthema}}\\[2cm]

    \begin{tabular}{rl}
        \textbf{Fach:} & \fach{} (\gerniveau) \\[0.3cm]
        \textbf{Klasse:} & \jahrgangsstufe \\[0.3cm]
        \textbf{Lehrwerk:} & \lehrwerk, \lehrwerkseite \\[0.3cm]
        \textbf{Datum:} & \unterrichtsdatum \\[0.3cm]
        \textbf{Zeit:} & \unterrichtszeit{} (Doppelstunde) \\[0.3cm]
        \textbf{Raum:} & \raum \\[0.3cm]
    \end{tabular}

    \vfill

    {\large Lehrkraft: \lehrkraft}\\[0.3cm]
    {\normalsize\textcolor{gray}{\ausbildungsstand}}\\[1cm]

\end{titlepage}

% ============================================================================
% INHALTSVERZEICHNIS
% ============================================================================
\tableofcontents
\newpage

% ============================================================================
% 1. BEDINGUNGSANALYSE
% ============================================================================
\section{Bedingungsanalyse}

\subsection{Lerngruppe}
\begin{tabular}{ll}
    Kursgröße: & \lerngruppengroesse{} SuS (\maennlich{} m / \weiblich{} w / \divers{} d) \\
    Jahrgangsstufe: & \jahrgangsstufe \\
    Sprachniveau: & \gerniveau{} (GER) \\
    Fremdsprache: & 2. Fremdsprache \\
\end{tabular}

\subsection{Differenzierung (intern)}

\textbf{WICHTIG:} Keine Sternchen auf Schülermaterial!

\begin{tabular}{|l|p{4cm}|p{4cm}|p{4cm}|}
\hline
& \textbf{[B] Basis} & \textbf{[S] Standard} & \textbf{[E] Erweiterung} \\
\hline
\textbf{Ziel} & Berufsreife & Mittlere Reife & Gymnasium \\
\hline
\textbf{Inhalt} & Konjugation mit Tabelle, Lückentext ausfüllen & Formen selbst bilden, einfache Sätze & Freie Urlaubsberichte verfassen \\
\hline
\textbf{Hilfen} & Konjugationstabelle, Signalwörter gegeben & Teilweise Hilfen & Nur Impulse \\
\hline
\end{tabular}

\subsection{Besonderheiten}
\begin{itemize}
    \item \textbf{LRS:} 2 SuS (Nachteilsausgleich: Zeitzugabe, vergrößerte AB)
    \item \textbf{DaZ:} 1 SuS (Wortschatz vorab klären)
    \item \textbf{Leistungsstark:} 4 SuS (Expertenrolle, Zusatzaufgaben)
\end{itemize}

% ============================================================================
% 2. SACHANALYSE
% ============================================================================
\section{Sachanalyse}

\subsection{Sprachliche Strukturen}

\textbf{Grammatik: Pretérito Indefinido regelmäßige -ar Verben}

Das Pretérito Indefinido ist eine Vergangenheitsform im Spanischen, die abgeschlossene Handlungen in der Vergangenheit beschreibt. Anders als das Pretérito Perfecto (Perfekt mit \textit{haber}) wird es für Ereignisse verwendet, die keinen Bezug zur Gegenwart haben.

\textbf{Konjugationsendungen -ar Verben:}
\begin{center}
\begin{tabular}{|l|l|l|}
\hline
\textbf{Person} & \textbf{Endung} & \textbf{Beispiel: hablar} \\
\hline
yo & -é & habl\textbf{é} \\
tú & -aste & habl\textbf{aste} \\
él/ella/usted & -ó & habl\textbf{ó} \\
nosotros/as & -amos & habl\textbf{amos} \\
vosotros/as & -asteis & habl\textbf{asteis} \\
ellos/ellas/ustedes & -aron & habl\textbf{aron} \\
\hline
\end{tabular}
\end{center}

\textbf{Behandelte Verben:}
\begin{itemize}
    \item \textbf{hablar} -- sprechen
    \item \textbf{viajar} -- reisen
    \item \textbf{nadar} -- schwimmen
    \item \textbf{tomar el sol} -- sonnenbaden
    \item \textbf{visitar} -- besuchen
\end{itemize}

\textbf{Signalwörter (palabras clave):}
\begin{itemize}
    \item \textbf{ayer} -- gestern
    \item \textbf{el año pasado} -- letztes Jahr
    \item \textbf{el verano pasado} -- letzten Sommer
    \item \textbf{la semana pasada} -- letzte Woche
\end{itemize}

\subsection{Kontrastive Betrachtung (Deutsch $\leftrightarrow$ Spanisch)}
\begin{itemize}
    \item \textbf{Positive Transfers:} Deutsche SuS kennen das Präteritum als abgeschlossene Vergangenheit
    \item \textbf{Interferenzen:} Im Deutschen wird oft Perfekt verwendet, wo Spanisch Indefinido nutzt
    \item \textbf{Besonderheit:} Akzente auf -é und -ó sind bedeutungsunterscheidend (hablo $\neq$ habló)
\end{itemize}

\subsection{Kulturelle Aspekte}
\begin{itemize}
    \item Urlaubserzählungen als typischer Kontext für Indefinido
    \item Beliebte Urlaubsziele spanischsprachiger Länder (Mallorca, Kanarische Inseln, Lateinamerika)
    \item Typische Urlaubsaktivitäten in hispanischen Kulturen
\end{itemize}

% ============================================================================
% 3. DIDAKTISCHE ANALYSE
% ============================================================================
\section{Didaktische Analyse}

\subsection{Stellung in der Sequenz}
Dies ist die \textbf{3. Doppelstunde} der Sequenz ``\reihenthema'' (4 Doppelstunden gesamt).

\textbf{Vorangegangen:}
\begin{itemize}
    \item 8a: Wortschatz Urlaub (Aktivitäten, Orte)
    \item 8b: Wegbeschreibung im Urlaubsort
\end{itemize}

\textbf{Folgend:}
\begin{itemize}
    \item 8d: Pretérito Indefinido -er/-ir Verben + Urlaubspräsentation
\end{itemize}

\subsection{Kompetenzziele (Rahmenplan MV)}
\begin{itemize}
    \item \textbf{Kommunikativ:} Vergangene Erlebnisse berichten (Schreiben, Sprechen)
    \item \textbf{Sprachlich:} Pretérito Indefinido Konjugation (-ar), Signalwörter
    \item \textbf{Interkulturell:} Urlaubskulturen vergleichen
    \item \textbf{Methodisch:} Induktives Erschließen von Grammatikregeln
\end{itemize}

\subsection{Legitimation}
\textbf{Rahmenplan Spanisch MV (Kl. 7-10):}
\begin{itemize}
    \item Kompetenzbereich: Sprachliche Mittel -- Grammatik
    \item Standard: ``Die SuS können über vergangene Ereignisse und Handlungen berichten.''
    \item Vergangenheitstempora als zentraler Bestandteil der A1/A2-Progression
\end{itemize}

% ============================================================================
% 4. STUNDENZIELE
% ============================================================================
\section{Stundenziele}

\subsection{Grobziel}
Die SuS erweitern ihre \textbf{kommunikative Kompetenz} im Bereich Schreiben und Sprechen, indem sie das Pretérito Indefinido regelmäßiger -ar Verben konjugieren und in einfachen Urlaubsberichten anwenden.

\subsection{Feinziele (operationalisiert)}
\begin{tabular}{|c|p{8cm}|c|c|}
\hline
\textbf{Nr.} & \textbf{Die SuS können...} & \textbf{AFB} & \textbf{Diff.} \\
\hline
FZ1 & die Konjugationsendungen des Pretérito Indefinido (-ar) korrekt wiedergeben & I & alle \\
\hline
FZ2 & die Verben hablar, viajar, nadar, tomar, visitar konjugieren & II & alle \\
\hline
FZ3 & Signalwörter (ayer, el año pasado, etc.) erkennen und zuordnen & I & alle \\
\hline
FZ4 & einfache Sätze über vergangene Urlaubserlebnisse bilden & II & [S]/[E] \\
\hline
FZ5 & einen kurzen zusammenhängenden Urlaubsbericht verfassen & III & [E] \\
\hline
\end{tabular}

% ============================================================================
% 5. METHODISCHE ANALYSE
% ============================================================================
\section{Methodische Analyse}

\subsection{Phasierung (Doppelstunde = 2 × 45 Min.)}
\begin{enumerate}
    \item \textbf{1. Stunde:}
    \begin{itemize}
        \item Einstieg: Bildimpuls Urlaubsfotos, Aktivierung Vorwissen
        \item Erarbeitung: Induktive Grammatikentdeckung (Textbeispiel)
        \item Systematisierung: Konjugationstabelle gemeinsam erstellen
        \item Übung 1: Chorisches Sprechen der Formen
        \item Übung 2: Zuordnungsübung (Signalwörter)
    \end{itemize}
    \item \textbf{2. Stunde:}
    \begin{itemize}
        \item Wiederholung: Konjugations-Quiz (Plenum)
        \item Übung 3: Lückentext (EA/PA)
        \item Anwendung: Sätze/Minitext über eigenen Urlaub schreiben
        \item Sicherung: Austausch in PA, Vorstellen ausgewählter Texte
        \item Abschluss: Zusammenfassung, Hausaufgabe
    \end{itemize}
\end{enumerate}

\subsection{Medien}
\begin{itemize}
    \item Tafel/Whiteboard (Konjugationstabelle)
    \item Lehrwerk: \lehrwerk, \lehrwerkseite
    \item Arbeitsblatt: AB-08c.pdf
    \item Lernpfad: LP-08c.html (digital, optional)
    \item Bildimpuls: Urlaubsfotos (Beamer/Poster)
\end{itemize}

% ============================================================================
% 6. VERLAUFSPLANUNG
% ============================================================================
\section{Verlaufsplanung}

\textbf{1. Stunde (45 Min.)}

\begin{longtable}{|p{1cm}|p{1.5cm}|p{4cm}|p{3cm}|p{2cm}|p{2cm}|}
\hline
\textbf{Zeit} & \textbf{Phase} & \textbf{Lehrerhandeln} & \textbf{Schülerhandeln} & \textbf{SF} & \textbf{Medien} \\
\hline
\endhead

0--5 & Einstieg & Zeigt Urlaubsfotos, fragt: ``¿Qué hicisteis en las vacaciones?'' & Vermutungen, bekannte Wörter nennen & Plenum & Beamer \\
\hline

5--12 & Input & Präsentiert Beispieltext (Urlaubsbericht), markiert Verbformen & Lesen, Verbformen identifizieren & Plenum & Tafel, Text \\
\hline

12--22 & Erar\-beitung & Fragt induktiv: ``Was fällt euch auf?'', leitet Regelfindung & Muster erkennen, Endungen formulieren & Plenum & Tafel \\
\hline

22--30 & System. & Erstellt Konjugationstabelle gemeinsam, erklärt Akzente & Ergänzen, Notieren & Plenum & Tafel, AB \\
\hline

30--38 & Übung 1 & Chorisches Sprechen anleiten: ``hablar -- hablé, hablaste...'' & Nachsprechen im Chor & Plenum & -- \\
\hline

38--45 & Übung 2 & Zuordnungsübung Signalwörter anleiten & Signalwörter zuordnen (AB Aufg. 1) & EA & AB \\
\hline
\end{longtable}

\vspace{1em}

\textbf{2. Stunde (45 Min.)}

\begin{longtable}{|p{1cm}|p{1.5cm}|p{4cm}|p{3cm}|p{2cm}|p{2cm}|}
\hline
\textbf{Zeit} & \textbf{Phase} & \textbf{Lehrerhandeln} & \textbf{Schülerhandeln} & \textbf{SF} & \textbf{Medien} \\
\hline
\endhead

0--8 & Wieder\-holung & Konjugations-Quiz: ``Wie heißt ,ich schwamm'?'' & Melden, Antworten & Plenum & Tafel \\
\hline

8--20 & Übung 3 & Lückentext erklären, differenzierte Hilfe geben & Lückentext ausfüllen (AB Aufg. 2--3) & EA/PA & AB \\
\hline

20--35 & Anwen\-dung & Schreibauftrag erteilen: ``Escribe sobre tus vacaciones'' & Eigene Sätze/Minitext schreiben (AB Aufg. 4) & EA & AB \\
\hline

35--42 & Sicherung & Austausch anleiten, ausgewählte Texte vorlesen lassen & Partneraustausch, Vorlesen & PA/Plenum & -- \\
\hline

42--45 & Abschluss & Zusammenfassung, HA: Lernpfad + Vokabeln & Notieren & Plenum & Tafel \\
\hline
\end{longtable}

\textbf{Legende:} EA = Einzelarbeit, PA = Partnerarbeit, SF = Sozialform

% ============================================================================
% 7. ERWARTETE SCHWIERIGKEITEN
% ============================================================================
\section{Erwartete Schwierigkeiten}

\begin{tabular}{|p{5cm}|p{8cm}|}
\hline
\textbf{Schwierigkeit} & \textbf{Reaktion} \\
\hline
Akzente auf -é und -ó vergessen & Bedeutungsunterschied betonen (hablo = ich spreche, habló = er sprach), visuelle Markierung an Tafel \\
\hline
Verwechslung Präsens/Indefinido & Signalwörter als ``Zeitmarker'' einführen, kontrastive Übung \\
\hline
Endung -amos identisch mit Präsens & Hinweis: Kontext und Signalwörter beachten \\
\hline
Schwache SuS bei freier Produktion & Satzbaukasten als Hilfe, Partnerarbeit mit stärkerem Partner \\
\hline
Zeitproblem & Puffer: Vorlesen kürzen, HA erweitern \\
\hline
\end{tabular}

% ============================================================================
% 8. HAUSAUFGABE
% ============================================================================
\section{Hausaufgabe}

\begin{itemize}
    \item Lernpfad LP-08c.html bearbeiten (digital)
    \item Vokabeln lernen: Signalwörter + Konjugation hablar
    \item Optional [E]: Eigenen Urlaubsbericht erweitern (5--8 Sätze)
\end{itemize}

% ============================================================================
% 9. ANLAGEN
% ============================================================================
\section{Anlagen}

\begin{enumerate}
    \item Arbeitsblatt (AB-08c.pdf)
    \item Musterlösung (ML-08c.pdf)
    \item Lehrerhinweise (LH-08c.pdf)
    \item Lernpfad (LP-08c.html)
    \item Beispieltext ``Mis vacaciones'' (Tafel/Folie)
    \item Urlaubsfotos (Beamer)
\end{enumerate}

% ============================================================================
% 10. DIDAKTIK-SCORE
% ============================================================================
\newpage
\section{Didaktik-Score}

\begin{scorebox}
\textbf{Bewertung der didaktischen Qualität nach 10 Dimensionen}\\
\textbf{Ziel-Score: $\geq$ 9.0 (Premium-Qualität)}

\vspace{1em}
\begin{tabular}{|c|l|c|c|p{4.5cm}|}
\hline
\textbf{\#} & \textbf{Dimension} & \textbf{Gew.} & \textbf{Score} & \textbf{Notiz} \\
\hline
1 & Differenzierung & 15\% & 9/10 & [B] mit Tabelle, [S] selbstständig, [E] freie Texte \\
\hline
2 & Taxonomie (AFB) & 10\% & 9/10 & I: 30\%, II: 45\%, III: 25\% \\
\hline
3 & Lernziel-Alignment & 10\% & 10/10 & Rahmenplan MV, A1-Progression \\
\hline
4 & Aktivierung & 10\% & 9/10 & Induktive Erarbeitung, eigene Texte \\
\hline
5 & Scaffolding & 10\% & 9/10 & Gestufte Hilfen, Satzbaukasten \\
\hline
6 & Feedback-Qualität & 10\% & 9/10 & Lernpfad mit Sofortfeedback \\
\hline
7 & Sprachliche Klarheit & 10\% & 9/10 & Altersgerecht Kl. 7 \\
\hline
8 & Motivation \& Relevanz & 10\% & 10/10 & Urlaubsthema, persönlicher Bezug \\
\hline
9 & Inklusion (UDL) & 10\% & 9/10 & Visuell, auditiv, Zeitzugabe LRS \\
\hline
10 & Praktikabilität & 5\% & 9/10 & 90 Min. realistisch, Materialien vorhanden \\
\hline
\multicolumn{3}{|r|}{\textbf{GESAMT:}} & \textbf{9.2/10} & \textcolor{successgreen}{\textbf{FREIGABE}} \\
\hline
\end{tabular}

\vspace{1em}
\textbf{Bewertungsschwellen:}
\begin{itemize}
    \item[\checkmark] \textcolor{successgreen}{$\geq$ 9.0: \textbf{FREIGABE} (Premium-Qualität)}
    \item[$\Box$] \textcolor{warningorange}{8.0--8.9: Iteration empfohlen}
    \item[$\Box$] 6.0--7.9: Überarbeitung erforderlich
    \item[$\Box$] $<$ 6.0: Grundlegende Neukonzeption
\end{itemize}
\end{scorebox}

\end{document}
